% AI-written, AI cross-reviewed; no human sign-off; not a claim that RH is proved
% RH Notes blueprint source excerpts for Z2.K6
% Verbatim excerpts in the notation of the parent blueprint. The labelled editions and source ranges appear below.
% See the unit YAML for the current proof status. Earlier-round passages retain their original wording.

% Start of atlas edition annotation. The following block is outside the original excerpt.
\paragraph{Status of named inputs in this reading.}
The citation identifies the in-repository class-group argument. The hand exposition invokes Artin reciprocity. Equivalence to the bibliographic formulation has no recorded comparison.
\noindent\nolinkurl{PositivityTheory.QuaternionShaBounds.ShaAsm_holds} --- in-repo constant (\path{Development/AutomorphicArithmetic/AutomorphicVarieties/ModularityArgument/ClassCounting/RingEquivalences/OrderOfOneAddPrimePowZmodPrimePow.lean:454--486}); not compared.\par
External theorem numbers in the original excerpt are not verified here unless a separate edition-specific source record documents the comparison.
% End of atlas edition annotation.

\section*{Z2.K6. Uniform layer count and the global kernel}

% Parent source SHA-256: 1af64922b98ed2310744411b1f18b25cc426c97220646321471836e1a3a409d6
% blueprint/ch_P.tex lines 1004 through 1045; lem:P-r3-k6
\begin{lemma}[Lifting the locally trivial characters (K6)]\label{lem:P-r3-k6}
Every character in the restricted localization kernel lifts to the
finite ideal class characters of Proposition~\ref{prop:P-k-count}
at every sufficiently large level. Their evaluation in $\mathcal H_q$
is injective.
\end{lemma}
\begin{proof}
By Lemma~\ref{lem:P-r3-k1}, restriction to the infinite layer is
injective. Represent a class over $M$ by a continuous cocycle $z$.
Its image is finite, and at each of the finitely many $w\mid\ell$
there is a $b_w\in A$ with $z(d)=(d-1)b_w$ on a chosen local group.
Choose $\nu_0\ge c$ so that $\ell^{\nu_0}$ kills the image and all
$b_w$. For every $\nu\ge\nu_0$ the group over $F=M(\mu_{\ell^\nu})$
acts trivially on $A[\ell^{\nu_0}]$. Thus $z$ restricts to a homomorphism
on this group. It vanishes on the chosen local groups above $\ell$,
and also on all conjugates there: for $h$ in the group over $F$,
\[
 z(dhd^{-1})=d z(h)+(1-dhd^{-1})z(d)=d z(h).
\]
Inertia away from $\ell$ already acts through the trivial group in the
unramified quotient. The resulting finite abelian character is unramified
at every place and split over $\ell$. Global Artin reciprocity identifies
it with an ideal class character, zero on primes over $\ell$. The
ordinary ideal class group applies since $F$ is totally complex. On
restriction back to the infinite layer the original value of $z$ is
recovered. These are the conditions in \nolinkurl{ClDual},
\nolinkurl{good_of}, \nolinkurl{ShaAsm_holds}
(\path{KummerTheory/Mk0IdMapCongr.lean:56--60},
\path{OrderOfOneAddPrimePowZmodPrimePow.lean:378--410,454--463}).
See \cite[Theorem 0.4 and the Artin map in Chapter V]{P-R3-CFT}.

Identify $A[q]$ additively with $\mathbf Z/q$. Frobenius at
$\gamma^{-1}P$ is conjugate to the transport of Frobenius at $P$.
The action of $\gamma^{-1}$ on the target is multiplication by
$(1+\ell^c)^E$, giving the required eigenvalue. Evaluation kills
principal divisors, and the prime ideals generate the ideal class
group, so two characters with the same evaluations agree. This proves
the assertion uniformly for all $\nu\ge\nu_0$ for the chosen class
(\nolinkurl{ClInj_holds}, \nolinkurl{sSha},
\path{OrderOfOneAddPrimePowZmodPrimePow.lean:265,490}).
\emph{Status: complete.}
\end{proof}

